\graphicspath{{content/chapters/4_methodology/figures/}}

\chapter{Methodology}%
\label{chp:methodology}

\section{System Overview}%
\label{sec:system_overview}

The overall architecture of the system consists of several components that work and interact together.
The main components include an intent classification module, a response generation module, and a conversational interface.


The user interacts with the system through the conversational interface by submitting a natural-language grammar question in English. 
The chatbot takes input in English as the students' proficiency in Latin may still be too limited to ask questions in Latin.
The user's input is processed by an intent classification model that identifies the grammatical category being asked about in the question.
The Latin word is extracted from the question and the system retrieves the relevant grammatical information from the Latin grammar knowledge base or a morphological analysis model.
A response is generated based on the identified intent and the retrieved grammatical information by using pre-defined response templates.
The response may combine both English and Latin to provide more exposure to the Latin language and encourage more interaction in Latin.
Finally, the chatbot returns the response to the user.


\section{Dataset Creation}%
\label{sec:dataset_creation}

\subsection{Latin Words Dataset}
\label{sec:latin_words_dataset}

% A dataset of Latin words and their grammatical features
% will be created by using the UniMorph dataset as the
% base dataset and using the Whitaker’s Words program
% to fill in and add additional grammar information to the
% entries.

A dataset of Latin words and their grammatical features
was created by using the \gls{unimorph} dataset as the
base dataset and using the Whitaker's Words program
to fill in and add additional grammar information to the
entries.

The \gls{unimorph} dataset was selected as it contains
many entries of individual Latin words, along with their
lemma and the grammatical features for how the word
was inflected. Having a large collection of words and
their corresponding grammatical features is needed to
construct the intent classification dataset, which will
contain questions about the grammar categories of specific words. While UniMorph is a good starting point
for data, it has some limitations. The first limitation is
that it does not include any information about the declension numbers of Latin nouns, nor does it include the
conjugation number of Latin verbs. These features are
important for Latin learners as they reveal the paradigms
of word inflections according to different grammatical
properties. Another limitation of \gls{unimorph} is that for
certain parts of speech, some grammatical properties
are left out. For example, gender is not included with
nouns, even though Latin nouns do have a grammatical
gender.

To address these limitations, Whitaker's Words will
be used to fill in the gaps. This program was chosen
since the output that it gives is already similar to the
format of the \gls{unimorph} entries, and it includes the
additional information that is missing. Since it is a
command line program, it can be easily interacted with
and controlled programmatically so that the extraction
of the needed information can be automated.

The result was a CSV file of many entries of Latin
word forms along with their grammatical features.

\subsection{Intent Classification Dataset}
\label{sec:intent_classification_dataset}

\textbf{question templates}

the templates are stored in a json file.
it consists of "placeholders" and "intents".
the placeholders are the variables that can be replaced with actual values to generate questions.
for example, WORD is a placeholder that can be replaced with any latin word from the dataset to generate a question.
other placeholders include NUMBER, TENSE, PERSON, etc.
these can be replaced with the corresponding values to generate questions.
these placeholders allow for the generation of a large number of questions from a limited set of templates, as they can be combined in various ways to create different questions.
the intents are the categories or labels that the questions belong to.
each intent represents a specific type of question that the system should be able to understand and respond to.
for example, an intent could be "conjugation" which would include questions related to verb conjugation, or "declension" which would include questions related to noun declension, etc.
each intent contains the list of target labels that the questions under that intent can be classified into.
and also the list of templates.
each template contains the text and the slots.
the text is the actual question template with the placeholders, and the slots are the corresponding placeholders that need to be filled in to generate a complete question.

%%%%%%% 
% TODO: UPDATE THIS SECTION

A structured dataset of grammar-related question templates 
and intents was created to be the training data for the 
intent classification model discussed in Section~\ref{sec:intent_classification_model}. 

The intent categories consist of the grammar feature 
the user is asking about for a specific word. For example, 
if the user is asking about what tense a verb is in, that 
question will be labelled as a tense query. The chosen intent 
categories are based on the different grammatical features 
covered in the dataset described in Section~\ref{sec:latin_words_dataset}. 
These categories should cover common questions learners may 
have about specific Latin words. Another category will be 
defined to cover out of scope inputs.

The dataset will be created by first creating example question 
templates for each intent category. The templates will include 
placeholders so that when generating the questions, the 
appropriate Latin words will be inserted in their place. 
The templates will be created by manually writing different 
example questions, as well as making use of large language 
models to help generate variations of different ways of how 
the same question can be phrased. The data augmentation 
techniques that will be used to increase the number of 
different question templates include paraphrasing of questions, 
both by replacing synonyms and varying the structure of the 
sentences. 

These templates will then be used to automatically generate 
many variations of questions to train the model, by taking a 
Latin word from the dataset described in the previous section 
and inserting it in the template. Along with each generated 
question, the corresponding intent will also be recorded in 
the dataset. This will create question and intent pairs for 
the classification task.


\subsection{Morphological Analysis Dataset}
\label{sec:morphological_analysis_dataset}

As explained in section~\ref{sec:fine_tuning_latin_bert_for_morphological_analysis}, two separate morphological analysers were created for this project, one trained on individual words and the other trained on sentences, and thus two separate datasets were used for training these morphological analysers.
The dataset used for training the word-level morphological analyser was the Latin Word Forms dataset discussed in section~\ref{sec:latin_words_dataset}, which contains individual Latin words along with their morphological features.

% 2 morphological analysers were created, trained on 2 separate datasets.
% one trained on individual words and the other trained on sentences.
% the dataset for the word-level morphological analyser was the Latin Word Forms dataset mentioned previously, which contains individual Latin words along with their morphological features.
for the sentence-level morphological analyser, existing Latin treebanks such as the PROIEL Treebank and the Perseus Latin Dependency Treebank were used, which contain annotated sentences with morphological information for each word.
these datasets do not contain complete morphological analyses for each word.
to address this, the missing morphological features were filled in automatically using existing tools, to create a more comprehensive dataset for training the morphological analyser.

\section{Intent Classification Model}%
\label{sec:intent_classification_model}

a bert based model was created to classify the grammatical intent of user questions.
the model was trained on the intent classification dataset created as described in section~\ref{sec:intent_classification_dataset}.
the model takes as input a natural language question and outputs the predicted intent category.

The model is fine-tuned using the dataset mentioned in section~\ref{sec:intent_classification_dataset}.

\section{Fine-tuning Latin BERT for Morphological Analysis}
\label{sec:fine_tuning_latin_bert_for_morphological_analysis}

% train a separate classifier for each morphological feature (e.g., tense, person, number, etc.) using the Latin BERT model.

% classifiers to be trained for the following morphological features
% - part of speech
% - tense
% - person
% - mood
% - voice
% - number
% - case
% - gender
% - aspect
% - declension
% - conjugation

% classifying a single word vs sequence labelling for a sentence
% single word may be more ambiguous, can have multiple analyses
% but current dataset is only single words
% think about this more

2 separate morphological analysers were created for this project, one trained on individual words and the other trained on sentences.
the word-level and sentence-level morphological analysers were both fine-tuned on the Latin BERT model, which is a pre-trained language model specifically designed for the Latin language.
the results of both morphological analysers were evaluated and compared to determine which approach is more effective for the task of morphological analysis in Latin.

\section{Response Generation}
\label{sec:response_generation}

template-based response generation

the response generation module uses a template-based approach to generate responses based on the identified intent and the retrieved grammatical information.
for each intent category, a set of response templates is defined, which contain placeholders for the relevant grammatical information.
when a response needs to be generated, the appropriate template is selected based on the identified intent, and the placeholders in the template are filled with the retrieved grammatical information to create a complete response.
the latin word from the question is extracted in order to get the relevant grammatical information for that word, which is then used to fill in the placeholders in the response template.

\section{Conversational Interface Design}
\label{sec:conversational_interface_design}

A web-based interface was developed for the chatbot, allowing users to interact with the system through a user-friendly interface.
The interface was designed to be simple and intuitive, with a text input field for users to submit

First there is a landing page where users are given a bit of background on the project and instructions on how to use the chatbot.
Then there is the main chat interface where users can submit their questions and receive responses from the chatbot.
The chat interface displays the conversation history, allowing users to see their previous questions and the chatbot's responses.
After chatting with the chatbot, users are directed to a feedback form where they can provide their feedback on the chatbot's performance and their overall experience.

